\section{Colors across initial region boundaries}
\label{sec:area_law_initial_region_interfaces}

Retain the initial region indices, closed birth regions $X_i^0$, and
colors $\lambda(i)$ of \cite[Section~11]{OpenAI2026AreaLaw}. The origin,
finite endpoint set, and valid residue choices are arbitrary, with
$C\ge 2$ and $k_0\ge 50{,}000{,}000$.

% Source: 10-geometry.tex, prop:two-families, lines 172--177,
% 212--218 and 299--323; geometry:nonadjacent, lines 193--198;
% geometry:primary-pieces, lines 246--249;
% revision adc7f1241b42e322a6451854ab7e4b4c146bf78a.
% This concerns the initial family before stars and repairs.

\begin{theorem}[Opposite colors along a shared segment]
    \label{thm:area_law_initial_interface_colors}
    \lean{TNLean.PEPS.AreaLaw.Geometry.initialRegionColor_ne_of_segment_subset_inter}
    \leanok
    \uses{def:area_law_initial_birth_regions,def:area_law_initial_region_colors}
    Let $i\ne j$ be distinct initial region indices. If $u,v\in\mathbb R^2$
    are distinct and their entire closed segment satisfies
    \begin{align}
        [u,v] & \subseteq X_i^0\cap X_j^0,
    \end{align}
    then $\lambda(i)\ne\lambda(j)$.
\end{theorem}
\begin{proof}\leanok
    \uses{thm:area_law_initial_run_color,
        thm:area_law_belt_run_primary_interface,
        thm:area_law_belt_runs_interface_colors,
        thm:area_law_belt_run_dummy_interface,
        thm:area_law_fan_run_geometric_contact,
        thm:area_law_local_layer_indices,
        thm:area_law_dummy_contact_index,
        thm:area_law_primary_diameter_separation}
    When the shared segment meets a belt run and either a primary region or
    the dummy region, the primary-contact and dummy-contact interface rules
    give opposite colors,
    since a run has the same color as each of its constituent triangles.
    The belt-run interface rule handles runs in
    distinct indexed cells. If two distinct runs lie in the same cell, the
    shared segment contains two distinct common points, and
    the same-cell geometric contact rule gives opposite colors.

    Let primary regions $i=(k,J)$ and $j=(h,K)$ with $k\le h$ share a point
    $x$. They lie in neighboring closed layers, so $h\in\{k,k+1\}$. If
    $h=k$, then $J\ne K$, and positive separation gives some $t>0$ with
    \begin{align}
        0=\|x-x\|_\infty & \ge t,
    \end{align}
    which is impossible. Hence $h=k+1$, and
    \begin{align}
        \lambda(j)=k+1\ne k=\lambda(i)\quad\text{in }\mathbb Z/2\mathbb Z.
    \end{align}
    A primary region $i=(k,J)$ meeting the dummy region lies in layer $k=k_0$, so
    \begin{align}
        \lambda(i)=k_0\ne k_0-1=\lambda(*)\quad\text{in }\mathbb Z/2\mathbb Z.
    \end{align}
    Two dummy indices cannot be distinct, and these cases exhaust all pairs.
\end{proof}

The hypothesis requires a nondegenerate segment in $X_i^0\cap X_j^0$;
isolated common corners of regions do not force opposite colors.
